\section[Gamma differential polynomials and exact finite-part reduction]{Gamma differential polynomials\\and exact finite-part reduction}
\label{gamma:sec}

The abstract reduction used in this section is the classical
functional-monomial decomposition of differential polynomials.
It appears explicitly in Guo--Regensburger--Rosenkranz
\cite[Proposition~4.10]{Guo2014}, with an earlier reduction rule in
Rosenkranz--Regensburger \cite[Lemma~5]{Rosenkranz2008}.
We give its short proof to specify the projection exactly.  The
special-function contribution is its application to the Gamma
hierarchy, the coefficient sectors in the next subsection, and the
endpoint and spectral formulas that follow.

\subsection{A canonical linear normal form}

Let $\mathscr R=\Q[u_0,u_1,u_2,\ldots]$, with derivation
$D u_r=u_{r+1}$ and weight $\wt(u_r)=r+1$.
The factor degree of $\prod u_r^{a_r}$ is $\sum a_r$.
A monomial is called functional if it is a power of $u_0$ (including
$1$), or its largest positive derivative index occurs at least twice.
Let $\mathscr F$ be their rational span.

\begin{theorem}[Classical normal form, with a fixed primitive]
\label{gamma:normal}
Every $P\in\mathscr R$ has a unique decomposition
\begin{equation}\label{gamma:decomp}
 P=\NN(P)+D\HH(P),\qquad \NN(P)\in\mathscr F,
 \qquad \HH(P)\text{ has zero constant term}.
\end{equation}
Both operators are rational and linear.  Reduction preserves factor
degree; $\NN$ preserves weight and $\HH$ lowers weight by one.
If $p(m)$ is the partition function, then
\begin{equation}\label{gamma:dimension}
 \dim_\Q\mathscr F_m=p(m)-p(m-1)\quad(m\ge2),
 \qquad \dim\mathscr F_0=\dim\mathscr F_1=1.
\end{equation}
\end{theorem}
\begin{proof}
If a monomial $M$ has a unique largest index $r\ge1$, replace that
factor $u_r$ by $u_{r-1}$, obtaining $\widetilde M$.
If $a$ is the multiplicity of $u_{r-1}$ in $\widetilde M$, then
\[
 M=\frac1aD\widetilde M+
      \left(M-\frac1aD\widetilde M\right).
\]
Every remaining monomial has largest derivative index at most $r-1$:
all differentiations of a factor $u_{r-1}$ give precisely the $a$
copies of $M$ that were removed.  Induction proves existence and
termination.  It also proves the grading assertions.

For uniqueness, a nonconstant polynomial $H$ of largest derivative
index $r$ has
$DH=u_{r+1}\partial H/\partial u_r+$ terms of order at most $r$.
Its terms of order $r+1$ have that largest index exactly once and
cannot belong to $\mathscr F$.  Hence $D\mathscr R\cap\mathscr F=0$,
and the kernel of $D$ consists of constants.  Fixing the constant of
$H$ proves uniqueness.  There are $p(m)$ monomials of weight $m$;
for $m\ge2$, $D$ is injective on the positive-weight space of weight
$m-1$.  This proves \eqref{gamma:dimension}.
\end{proof}

The quotient $\mathscr R/D\mathscr R$ in this theorem is a
\emph{vector-space quotient}.  For example $\NN(u_1)=0$, but
$\NN(u_1^2)=u_1^2$.  Thus $\NN$ is not a ring homomorphism.
No multiplication of congruences modulo $D\mathscr R$ is permitted
without a separate justification.

Now put $u_r=\psi^{(r)}(x)$.  At zero the convergent Laurent expansion is
\begin{equation}\label{gamma:local}
 \psi(x)=-\frac1x-\gamma+
       \sum_{j\ge1}(-1)^{j+1}\zeta(j+1)x^j,
       \qquad |x|<1.
\end{equation}
It follows from the standard log-Gamma series \cite[\S5.7]{DLMF};
all higher local expansions follow by differentiation.
Define the explicit endpoint functional
\begin{equation}\label{gamma:Delta}
 \Delta H=H\bigl(\psi(1),\psi'(1),\ldots\bigr)
             -[x^0]H\bigl(\psi(x),\psi'(x),\ldots\bigr).
\end{equation}
Then \eqref{scope:boundary} gives
\begin{equation}\label{gamma:fp-normal}
 \FP\int_0^1P(\psi,\psi',\ldots)\dd x
 =\FP\int_0^1\NN(P)(\psi,\psi',\ldots)\dd x+\Delta\HH(P).
\end{equation}
This is a finite algebraic recipe: for a monomial of weight $w$, write
it as $x^{-w}\prod_r(x^{r+1}\psi^{(r)}(x))^{a_r}$ and take its
coefficient of $x^w$.  Only terms through degree $w$ in these ordinary
series are needed.  The resulting endpoint belongs to
$\Q[\gamma,\zeta(2),\zeta(3),\ldots]$.

\subsection{The inverse-Gamma hierarchy and its balance identity}

Define
\begin{equation}\label{gamma:A}
 \A_m(x)=(-1)^m\Gamma(x)\frac{d^m}{dx^m}\frac1{\Gamma(x)},
 \qquad \A_0=1.
\end{equation}
As universal differential polynomials,
\begin{equation}\label{gamma:recurrence}
 \A_{m+1}=u_0\A_m-D\A_m,
 \qquad
 \sum_{m\ge0}\A_m\frac{t^m}{m!}
 =\exp\left(\sum_{r\ge0}\frac{(-1)^ru_rt^{r+1}}{(r+1)!}\right).
\end{equation}
The generating function is simply $\Gamma(x)/\Gamma(x-t)$.
Writing $B_j=\Gamma^{(j)}(x)/\Gamma(x)$ gives the exact balance
identity
\begin{equation}\label{gamma:balanced}
 \A_m=B_r\A_{m-r}
       -D\sum_{j=0}^{r-1}B_j\A_{m-1-j},\qquad 0\le r\le m.
\end{equation}
Indeed $DB_j=B_{j+1}-u_0B_j$ and
$D\A_k=u_0\A_k-\A_{k+1}$, so differentiating the sum telescopes.
Also $B_m(u)=(-1)^m\A_m(-u)$.  Taking $r=m$ in
\eqref{gamma:balanced} and then the normal form proves
\begin{equation}\label{gamma:parity}
 \NN(\A_m)(u)=(-1)^m\NN(\A_m)(-u).
\end{equation}
Only factor degrees congruent to $m$ modulo two survive.

The first reductions illustrate what must still be integrated:
\begin{align}
 \NN(\A_2)&=u_0^2,&
 \NN(\A_3)&=u_0^3,\notag\\
 \NN(\A_4)&=u_0^4-u_1^2,&
 \NN(\A_5)&=u_0^5-5u_0u_1^2,\notag\\
 \NN(\A_6)&=u_0^6-15u_0^2u_1^2+u_2^2,\label{gamma:firstforms}\\
 \NN(\A_7)&=u_0^7-35u_0^3u_1^2+7u_0u_2^2,\notag\\
 \NN(\A_8)&=u_0^8-70u_0^4u_1^2+28u_0^2u_2^2-7u_1^4-u_3^2.
 \notag
\end{align}
The full differential identities, including primitives, are exact
polynomial identities.  The accompanying code supplies them through
weight ten; the same finite algorithm applies at every weight.

\subsection{The inherited scalar bridge, with its normalization proved}

Let
\begin{equation}\label{gamma:c}
 \frac1{\Gamma(1+u)}=\sum_{j\ge0}c_j u^j,
 \qquad c_0=1,\quad c_1=\gamma.
\end{equation}
The source reports establish the all-depth Gamma-quotient bridge
\cite[\S3]{SourceCanonical}.  We reproduce a proof here because its
Laurent sign and scalar normalization enter the sixth moment.
Define
\begin{equation}\label{gamma:C}
 \mathcal C_m(x)=m![t^m]
       \frac{\Gamma(1-t)\Gamma(x)}{\Gamma(x-t)}.
\end{equation}
For $m\ge2$ it extends continuously to $[0,1]$, and
\begin{equation}\label{gamma:Cintegral}
 \int_0^1\mathcal C_m(x)\dd x=m!(c_m-\lambda_{m-1}).
\end{equation}

Here is a direct Mellin normalization.  Set
$L(v)=-\log(1-e^{-v})$ and
\[
 V_d(v)=L(v)^d-\boldsymbol1_{(0,1)}(v)(-\log v)^d.
\]
The generating identity for strict unit-index harmonic sums and
integration by parts give, initially for $\re u>0$,
\begin{equation}\label{gamma:mellinheight}
 \zeta(1+u,\{1\}^{d-1})
 =\frac1{\Gamma(1+u)u^d}
  +\frac{u}{d!\Gamma(1+u)}\int_0^\infty v^{u-1}V_d(v)\dd v.
\end{equation}
For clarity, the unsplit integral is
$[\Gamma(u)d!]^{-1}\int_0^\infty v^{u-1}L(v)^d\dd v$.
It follows by using
\[
 \sum_{n\ge1}H_{n-1}(\{1\}^{d-1})e^{-nv}
 =\frac{e^{-v}L(v)^{d-1}}{(d-1)!(1-e^{-v})}
 =-\frac1{d!}\frac d{dv}L(v)^d.
\]
The subtracted integral in \eqref{gamma:mellinheight} is holomorphic
near $u=0$: near zero $V_d(v)=O(v(1+|\log v|)^{d-1})$, and at
infinity it decays exponentially.  Consequently
\begin{equation}\label{gamma:lambda-integral}
 \lambda_d=c_{d+1}+\frac1{d!}\int_0^\infty\frac{V_d(v)}v\dd v.
\end{equation}

For the other side put
$F(x,t)=\Gamma(1+t)\Gamma(x)/\Gamma(x+t)$ and
\[
 \Phi(t)=\int_0^1\{F(x,t)-1+t(\psi(x)+\gamma)\}\dd x.
\]
The residue $t/x$ of $F$ cancels the residue $-t/x$ of the last
term.  The remaining integrand is jointly holomorphic in $t$ near
zero and continuous in $x\in[0,1]$.
Euler's beta integral, combined with the digamma integral before
interchanging integrations, yields
\begin{equation}\label{gamma:Phi}
 \Phi(t)=t\int_0^\infty
 \frac{(1-e^{-v})^t-1-\boldsymbol1_{(0,1)}(v)(v^t-1)}v\dd v.
\end{equation}
To verify the subtraction constant, for $\re t>0$ integration in
$x$ first gives
\[
 t\int_0^\infty
 \left\{\frac{(1-e^{-v})^t-1}{v}
              +\frac{e^{-v}}{1-e^{-v}}\right\}\dd v-1.
\]
With lower cutoff $\delta$, the second integral is
$-\log(1-e^{-\delta})$ and
$\int_\delta^1(v^t-1)\dd v/v=(1-\delta^t)/t+\log\delta$.
Their limits give exactly \eqref{gamma:Phi}.
The latter integral continues holomorphically to a neighborhood of
zero by its displayed local subtraction.  Comparing the coefficient
of $t^m$ on its two sides gives
$\int\mathcal C_m=-m\int V_{m-1}(v)\dd v/v$, proving
\eqref{gamma:Cintegral}.  This proof also justifies all coefficient
extractions as ordinary integrals.

\begin{theorem}[All inverse-Gamma finite parts]
\label{gamma:scalar}
For $m\ge2$ put $\mathcal L_m=\FP\int_0^1\A_m(x)\dd x$.
Then
\begin{equation}\label{gamma:L}
 \boxed{\displaystyle
 \mathcal L_m=K_m-m!\sum_{d=1}^{m-1}
       (-1)^{m-1-d}c_{m-1-d}\lambda_d,\qquad
 K_{2n}=\frac{(-1)^n\pi^{2n}}{2n+1},\quad K_{2n+1}=0.}
\end{equation}
This is the scalar bridge already implicit in the incoming all-order
Gamma formulas; the normal-form and spectral applications below
are the extensions developed here.
\end{theorem}
\begin{proof}
The coefficientwise finite part of the normalized quotient is
\[
 1+\gamma t+\sum_{m\ge2}(c_m-\lambda_{m-1})t^m
 =\frac1{\Gamma(1+t)}-\sum_{d\ge1}\lambda_d t^{d+1}.
\]
Here $\FP\int\mathcal C_1=\gamma$, since
$\mathcal C_1=\psi+\gamma$ and $\FP\int\psi=0$.
Multiply by $1/\Gamma(1-t)=\sum(-1)^jc_jt^j$.
Gamma reflection turns the scalar product into
$\sin(\pi t)/(\pi t)$.  Extracting $m![t^m]$ proves
\eqref{gamma:L}.
\end{proof}

If
$\NN(\A_m)=u_0^m+R_m$ and $H_m=\HH(\A_m)$, the immediate
all-order moment identity is
\begin{equation}\label{gamma:allmoments}
 \boxed{\displaystyle
 Q_m:=\FP\int_0^1\psi(x)^m\dd x
   =\mathcal L_m-\Delta H_m
       -\FP\int_0^1R_m(\psi,\psi',\ldots)\dd x.}
\end{equation}
Every monomial of $R_m$ contains a derivative.  Section~\ref{spectral:sec}
will evaluate that remainder by convergent harmonic coordinates and
provide its explicit antiderivative, so \eqref{gamma:allmoments} is
part of a terminating exact procedure rather than an instruction to
leave mixed integrals unnamed.
